% ROZDZIAŁ 19 — zastosowania pochodnej Radona--Nikodyma
\chapter{Zastosowania twierdzenia Radona--Nikodyma}
\label{ch:rn-applications-19}

\section{Gęstość jako opis zmiany miary}
Twierdzenie Radona--Nikodyma z rozdziału~\ref{ch:radon-nikodym-18}
zamienia porównywanie miar w rachunek funkcyjny. Nie wystarczy jednak
bezrefleksyjnie dzielić gęstości: przy każdej zmianie miary należy
wskazać, względem której miary zachodzi równość prawie wszędzie,
a także sprawdzić, czy wyrażenia są całkowalne.

\begin{twierdzenie}[Całkowanie po zmianie miary]
\label{thm:change-measure-19}
Niech $(X,\mathcal F)$ będzie przestrzenią mierzalną,
$\mu,\nu$ nieujemnymi miarami $\sigma$-skończonymi na niej,
$\nu\ll\mu$, a $f=\mathrm d\nu/\mathrm d\mu$.
Dla każdej mierzalnej funkcji $h:X\to[0,\infty]$ zachodzi
\begin{equation}\label{eq:change-measure-19}
 \int_X h(x)\,\mathrm d\nu(x)
 =\int_X h(x)f(x)\,\mathrm d\mu(x).
\end{equation}
Jeśli $h:X\to\mathbb C$ jest mierzalna, to $h\in L^1(\nu)$
wtedy i tylko wtedy, gdy $hf\in L^1(\mu)$; wówczas
\eqref{eq:change-measure-19} pozostaje prawdziwy.
\end{twierdzenie}
\begin{proof}
Dla $h=\chi_A$ równość jest definicją pochodnej.
Przez liniowość otrzymujemy ją dla nieujemnych funkcji prostych,
a następnie przez aproksymację monotoniczną i twierdzenie Beppo
Leviego dla każdej $h\ge0$. Zastosowanie wyniku do $|h|$
daje równoważność całkowalności. Rozłożenie funkcji zespolonej
na dodatnie i ujemne części jej części rzeczywistej i urojonej
kończy dowód.
\end{proof}

\begin{przyklad}[Ważenie miary]
Niech $\mu=\lambda_1|_{(0,1)}$ i $f(x)=2x$.
Wówczas $\nu(E)=\int_E 2x\,\mathrm d\lambda_1(x)$ jest
miarą probabilistyczną. Dla $h(x)=x^2$ otrzymujemy
$\int_{(0,1)}h\,\mathrm d\nu(x)
=\int_{(0,1)}2x^3\,\mathrm d\lambda_1(x)=1/2$.
Z kolei $\mathrm d\mu/\mathrm d\nu=(2x)^{-1}$
$\nu$-prawie wszędzie. Mimo nieograniczoności tej pochodnej
przy zerze zachodzi
$\int_{(0,1)}(2x)^{-1}\,\mathrm d\nu(x)=1$.
\end{przyklad}

\section{Miary obrazowe oraz gęstości obserwacji}
Dla mierzalnego $T:X\to Y$ i miary $\mu$ na $(X,\mathcal F)$
przypominamy miarę obrazową
$T_*\mu(B)=\mu(T^{-1}(B))$, $B\in\mathcal G$,
gdzie $(Y,\mathcal G)$ jest przestrzenią mierzalną.
Nie wymagamy, aby $T$ była odwracalna.

\begin{twierdzenie}[Bezwzględna ciągłość zachowana przez obraz]
\label{thm:image-ac-19}
Niech $T:(X,\mathcal F)\to(Y,\mathcal G)$ będzie mierzalne,
a $\mu,\nu$ miarami nieujemnymi na $\mathcal F$.
Jeżeli $\nu\ll\mu$, to $T_*\nu\ll T_*\mu$.
Jeżeli ponadto $\mu,\nu$ są miarami probabilistycznymi,
$f=\mathrm d\nu/\mathrm d\mu$ i
$g=\mathrm d(T_*\nu)/\mathrm d(T_*\mu)$,
to $g\circ T$ jest wersją warunkowej wartości oczekiwanej
$\mathbb E_\mu[f\mid\sigma(T)]$, gdzie
$\sigma(T)=\{T^{-1}(B):B\in\mathcal G\}$.
\end{twierdzenie}
\begin{proof}
Jeżeli $(T_*\mu)(B)=0$, to $\mu(T^{-1}(B))=0$,
zatem $\nu(T^{-1}(B))=0$, co jest równoważne
$(T_*\nu)(B)=0$. Obie miary obrazowe są skończone,
więc pochodna $g$ istnieje. Dla każdego $B\in\mathcal G$
z twierdzenia o całkowaniu względem miary obrazowej
(definicję miary obrazowej z rozdziału~\ref{ch:struktura-przestrzeni})
uzyskujemy
\[
\int_{T^{-1}(B)}g(T(x))\,\mathrm d\mu(x)
=\int_Bg(y)\,\mathrm d(T_*\mu)(y)
=(T_*\nu)(B)
=\int_{T^{-1}(B)}f(x)\,\mathrm d\mu(x).
\]
Funkcja $g\circ T$ jest $\sigma(T)$-mierzalna. Jest również
całkowalna, ponieważ
\[
\int_X|g(T(x))|\,\mathrm d\mu(x)
=\int_Y|g(y)|\,\mathrm d(T_*\mu)(y)<\infty.
\]
Spełnia więc warunek charakteryzujący warunkową wartość oczekiwaną
z definicji~\ref{def:conditional-exp-18}.
\end{proof}

\begin{uwaga}[Kiedy informacja znika?]
W ogólności $g(T(x))\neq f(x)$. Po zastosowaniu
odwzorowania $T$ pozostaje tylko średnia warunkowa $f$
względem informacji dostępnej poprzez $T$. Jeśli $T$
jest stałe, a $\mu,\nu$ są miarami probabilistycznymi,
to $T_*\mu=T_*\nu$, choć $\mu$ i $\nu$ mogą być różne.
Nie wolno też z samego $\nu\ll\mu$ wnioskować, że dowolne
odwzorowanie $T$ zachowuje osobliwość miar.
\end{uwaga}

\section{Wariacja całkowita a testowanie różnic między rozkładami}
Dla miar probabilistycznych $P,Q$ na $(X,\mathcal F)$
wprowadzamy odległość
\begin{equation}\label{eq:tv-prob-19}
d_{\mathrm{TV}}(P,Q):=\sup_{A\in\mathcal F}|P(A)-Q(A)|.
\end{equation}
Należy odróżnić ją od normy wariacji całkowitej miary ze znakiem
$P-Q$, oznaczonej w rozdziale~\ref{ch:miary-ze-znakiem} przez
$\|P-Q\|_{\mathrm{TV}}=|P-Q|(X)$.

\begin{twierdzenie}[Formuła wariacyjna dla dwóch prawdopodobieństw]
\label{thm:total-variation-formula-19}
Niech $P,Q$ będą miarami probabilistycznymi na $(X,\mathcal F)$.
Niech $\lambda=P+Q$, $p=\mathrm dP/\mathrm d\lambda$
i $q=\mathrm dQ/\mathrm d\lambda$. Wówczas
\begin{equation}\label{eq:tv-formula-19}
 d_{\mathrm{TV}}(P,Q)
 =\frac12\int_X|p(x)-q(x)|\,\mathrm d\lambda(x)
 =\frac12\|P-Q\|_{\mathrm{TV}}
 =1-\int_X\min(p,q)\,\mathrm d\lambda.
\end{equation}
Supremum z \eqref{eq:tv-prob-19} jest osiągane na zbiorze
$A_+=\{p>q\}$. W szczególności $0\le d_{\mathrm{TV}}(P,Q)\le1$,
a równość $d_{\mathrm{TV}}(P,Q)=1$ zachodzi dokładnie wtedy,
gdy $P\perp Q$.
\end{twierdzenie}
\begin{proof}
Miara ze znakiem $\eta=P-Q$ ma gęstość $p-q$
względem $\lambda$. Zatem
$\eta(A)=\int_A(p-q)\,\mathrm d\lambda$ i
$\eta(X)=0$. Rozkład na części dodatnią i ujemną
pokazuje, że $\eta^+(X)=\eta^-(X)=|\eta|(X)/2$.
Dla dowolnego $A$ mamy
$-\eta^-(X)\le\eta(A)\le\eta^+(X)$,
a $\eta(A_+)=\eta^+(X)$; stąd pierwsze dwie równości.
Ponieważ $\int_X(p+q)\,\mathrm d\lambda=2$ i punktowo
$|p-q|=p+q-2\min(p,q)$, dostajemy ostatnią równość.
Jeżeli $d_{\mathrm{TV}}=1$, to $\min(p,q)=0$
$\lambda$-prawie wszędzie, czyli $p q=0$ prawie wszędzie.
Zbiory $\{p>0\}$ i $\{q>0\}$ są wówczas rozłączne
z dokładnością do zbioru $\lambda$-zerowego i skupiają
odpowiednio $P$ i $Q$; stąd $P\perp Q$.
Implikacja odwrotna wynika bezpośrednio z definicji osobliwości.
\end{proof}

\begin{wniosek}[Przetwarzanie danych nie zwiększa wariacji całkowitej]
\label{cor:tv-data-processing-19}
W założeniach twierdzenia~\ref{thm:image-ac-19},
bez wymogu $\nu\ll\mu$, jeśli $P,Q$ są prawdopodobieństwami,
to
$ d_{\mathrm{TV}}(T_*P,T_*Q)\le d_{\mathrm{TV}}(P,Q)$.
\end{wniosek}
\begin{proof}
Dla dowolnego $B\in\mathcal G$ zachodzi
$|(T_*P)(B)-(T_*Q)(B)|
=|P(T^{-1}B)-Q(T^{-1}B)|\le d_{\mathrm{TV}}(P,Q)$.
Bierzemy supremum po $B$.
\end{proof}

\begin{przyklad}[Zbieżność słaba nie wymusza zbieżności w wariacji]
Na $\mathbb R$ rozważmy $P_n=\delta_{1/n}$ i $P=\delta_0$.
Dla każdej funkcji ciągłej i ograniczonej $h$
zachodzi $\int_{\mathbb R}h\,\mathrm dP_n\to
\int_{\mathbb R}h\,\mathrm dP$, co jest zbieżnością wąską.
Jednak $d_{\mathrm{TV}}(P_n,P)=1$ dla każdego $n$,
ponieważ $P_n(\{1/n\})=1$ i $P(\{1/n\})=0$.
Precyzyjną teorię zbieżności miar rozwiniemy w rozdziale~22.
\end{przyklad}

\section{Lemat Neymana--Pearsona: optymalność testu progowego}
Przedstawimy rezultat stanowiący użyteczne, a zarazem nieintuicyjne
zastosowanie pochodnej Radona--Nikodyma. Z dwóch prostych hipotez
można skonstruować test optymalny bez żadnych założeń o istnieniu
gęstości względem miary Lebesgue'a.

\begin{definicja}[Test losowy i jego moc]
\label{def:test-19}
Niech $P,Q$ będą prawdopodobieństwami na $(X,\mathcal F)$.
Testem losowym jest mierzalna funkcja $\varphi:X\to[0,1]$.
Liczbę $\int_X\varphi\,\mathrm dP$ nazywamy
prawdopodobieństwem odrzucenia hipotezy $P$ pod $P$,
a $\int_X\varphi\,\mathrm dQ$ jego mocą pod $Q$.
Test ma poziom nieprzekraczający $\alpha\in[0,1]$,
gdy pierwsza z tych liczb nie przekracza $\alpha$.
\end{definicja}

\begin{twierdzenie}[Lemat Neymana--Pearsona]
\label{thm:neyman-pearson-19}
Niech $P,Q$ będą prawdopodobieństwami na $(X,\mathcal F)$,
$\lambda=P+Q$, $p=\mathrm dP/\mathrm d\lambda$
i $q=\mathrm dQ/\mathrm d\lambda$.
Dla $c\ge0$ i $\gamma\in[0,1]$ ustalmy test
\[
\varphi_*=\chi_{\{q>cp\}}+
 \gamma\chi_{\{q=cp\}}.
\]
Niech $\alpha=\int_X\varphi_*\,\mathrm dP$.
Wtedy dla każdego testu $\varphi:X\to[0,1]$ o poziomie
nieprzekraczającym $\alpha$ zachodzi
\[
 \int_X\varphi\,\mathrm dQ\le
 \int_X\varphi_*\,\mathrm dQ.
\]
Dla $c>0$ dowód nie wymaga założenia, że
$P$ i $Q$ są wzajemnie bezwzględnie ciągłe.
W razie $c=0$ twierdzenie pozostaje prawdziwe;
nie wymaga istnienia gęstości ani $P$, ani $Q$ względem
żadnej wcześniej zadanej miary odniesienia.
\end{twierdzenie}
\begin{proof}
Na zbiorze $\{q>cp\}$ mamy $\varphi_*=1$, więc
$\varphi-\varphi_*\le0$; na $\{q<cp\}$ mamy
$\varphi_*=0$, więc $\varphi-\varphi_*\ge0$.
Na zbiorze $\{q=cp\}$ iloczyn jest zerem. W każdym
przypadku zatem
$(\varphi-\varphi_*)(q-cp)\le0$ punktowo.
Całkujemy względem $\lambda$ i dostajemy
\[
\int_X(\varphi-\varphi_*)\,\mathrm dQ
\le c\int_X(\varphi-\varphi_*)\,\mathrm dP\le0,
\]
gdzie ostatnia nierówność wynika z poziomu testu
i nieujemności $c$. To dokładnie teza.
\end{proof}

\begin{uwaga}[Rola losowania na granicy]
Wprowadzony parametr $\gamma$ pozwala kontrolować poziom
testu, gdy zbiór $\{q=cp\}$ ma dodatnią miarę $P$.
W ogólności nie każdy poziom $\alpha$ da się otrzymać
bez losowania. Na zbiorze $\{p=0,q>0\}$ test progowy
zawsze ma wartość $1$: odrzucenie $P$ nic tam nie kosztuje.
\end{uwaga}

\begin{wniosek}[Konstrukcja testu optymalnego o zadanym poziomie]
\label{cor:NP-exact-alpha-19}
Niech $P,Q,\lambda,p,q$ mają znaczenie z twierdzenia
\ref{thm:neyman-pearson-19}. Dla każdego
$\alpha\in(0,1)$ istnieją $c\ge0$ i
$\gamma\in[0,1]$, dla których test
$\varphi_*=\chi_{\{q>cp\}}+\gamma\chi_{\{q=cp\}}$
ma poziom dokładnie $\alpha$, a spośród wszystkich
testów poziomu co najwyżej $\alpha$ ma największą moc.
Dla $\alpha=0$ testem optymalnym jest
$\chi_{\{p=0,\,q>0\}}$, a dla $\alpha=1$
można przyjąć test stały równy $1$.
\end{wniosek}
\begin{proof}
Niech $L=q/p$ na $\{p>0\}$ i $L=0$ na $\{p=0\}$.
Ponieważ $P(\{p=0\})=0$, mamy $L<\infty$
$P$-prawie wszędzie. Dla $0<\alpha<1$ kładziemy
\[
 c=\inf\{t\ge0:P(L>t)\le\alpha\}.
\]
Z nierówności $P(L>t)\downarrow0$ przy $t\uparrow\infty$
wynika, że zbiór pod infimum jest niepusty, a $c$ jest
skończone. Dla $t_n\downarrow c$ i $t_n>c$ z ciągłości
miary z dołu otrzymujemy $P(L>c)\le\alpha$.
Jeśli $c>0$, to z definicji infimum
$P(L>t)>\alpha$ dla $0\le t<c$;
przechodząc z $t\uparrow c$ i używając ciągłości miary
z góry, dostajemy $P(L\ge c)\ge\alpha$.
Gdy $c=0$, ostatnia nierówność wynika z
$P(L\ge0)=1$.
W konsekwencji
\[
 P(L>c)\le\alpha\le P(L\ge c).
\]
Jeśli $P(L=c)>0$, przyjmujemy
$\gamma=(\alpha-P(L>c))/P(L=c)\in[0,1]$;
w przeciwnym przypadku nierówności dają
$P(L>c)=\alpha$ i wybieramy $\gamma=0$.
Zbiory $\{L>c\}$, $\{L=c\}$ zgadzają się
$P$-prawie wszędzie z $\{q>cp\}$ i $\{q=cp\}$,
więc $\int_X\varphi_*(x)\,\mathrm dP(x)=\alpha$.
Optymalność wynika z twierdzenia
\ref{thm:neyman-pearson-19}.

Dla $\alpha=0$ dopuszczalny test $\varphi$ jest zerowy
$P$-prawie wszędzie, a więc $\varphi=0$
$\lambda$-prawie wszędzie na $\{p>0\}$.
Stąd jego moc nie przekracza
$Q(\{p=0\})$, osiąganej przez wskazany test.
Przypadek $\alpha=1$ jest natychmiastowy.
\end{proof}

\begin{przyklad}[Dlaczego losowanie na granicy testu bywa konieczne?]
\label{ex:np-randomization-19}
Niech $X=\{a,b\}$, $\mathcal F=\mathcal P(X)$ oraz
$P(\{a\})=7/10$, $P(\{b\})=3/10$,
$Q(\{a\})=1/5$ i $Q(\{b\})=4/5$.
Stosunek wiarygodności $q/p$ dla $\lambda=P+Q$
przyjmuje wartości $2/7$ w punkcie $a$ i $8/3$
w punkcie $b$. Dla poziomu $\alpha=3/20$
przyjmujemy $c=8/3$ oraz $\gamma=1/2$.
Wówczas $\{q>cp\}=\varnothing$, a $\{q=cp\}=\{b\}$,
zatem test progowy ma postać
$\varphi_*=\frac12\chi_{\{b\}}$. Jego poziom jest
równy $\int_X\varphi_*\,\mathrm dP=3/20$, a moc
$\int_X\varphi_*\,\mathrm dQ=2/5$.

Każdy test \emph{nielosowy} przyjmujący tylko wartości
$0$ i $1$, którego poziom nie przekracza $3/20$,
musi być identycznie zerowy: odrzucenie hipotezy
w którymkolwiek punkcie kosztuje co najmniej
$P(\{b\})=3/10$. Ma więc moc równą $0$.
Losowanie na $\{q=cp\}$ rzeczywiście zwiększa
moc dopuszczalnego testu.
\end{przyklad}

\section{Wzór Bayesa jako twierdzenie o gęstościach}
Wynik sformułujemy bez założenia, że zmienna obserwowana
posiada klasyczną gęstość na $\mathbb R^d$.

\begin{twierdzenie}[Bayes dla skończonej przestrzeni parametrów]
\label{thm:bayes-finite-19}
Niech $m\ge1$, $\pi_i>0$, $\sum_{i=1}^m\pi_i=1$.
Niech $(Y,\mathcal G)$ będzie przestrzenią mierzalną,
a $P_1,\ldots,P_m$ miarami probabilistycznymi na $Y$.
Na przestrzeni $\{1,\ldots,m\}\times Y$ definiujemy
prawdopodobieństwo wspólne przez
$\mathbb P(\{i\}\times B)=\pi_iP_i(B)$.
Niech $\lambda=\sum_iP_i$, $p_i=\mathrm dP_i/\mathrm d\lambda$,
a $r=\sum_i\pi_i p_i$. Dla $r(y)>0$ określmy
\[
 \beta_i(y)=\frac{\pi_i p_i(y)}{r(y)}.
\]
Na zbiorze $\{r=0\}$ możemy przyjąć dowolne mierzalne
wartości $\beta_i\in[0,1]$. Wtedy dla obserwacji
$Y(i,y)=y$ i etykiety $\Theta(i,y)=i$ zachodzi
\begin{equation}\label{eq:bayes-finite-19}
 \beta_i(Y)=\mathbb P(\Theta=i\mid\sigma(Y))
 \quad\mathbb P\text{-prawie wszędzie}.
\end{equation}
Ponadto $\sum_i\beta_i=1$ $P_Y$-prawie wszędzie,
gdzie $P_Y=\sum_i\pi_iP_i$ jest prawem obserwacji.
\end{twierdzenie}
\begin{proof}
Dla dowolnego $B\in\mathcal G$ definicja rozkładu
wspólnego daje
$\mathbb P(\Theta=i,\,Y\in B)=\pi_iP_i(B)
=\int_B\pi_ip_i(y)\,\mathrm d\lambda(y)$.
Tymczasem $P_Y(B)=\int_Br(y)\,\mathrm d\lambda(y)$,
więc stosując twierdzenie~\ref{thm:change-measure-19}
uzyskujemy
\[
\int_{Y^{-1}(B)}\beta_i(Y)\,\mathrm d\mathbb P
=\int_B\beta_i(y)r(y)\,\mathrm d\lambda(y)
=\pi_iP_i(B).
\]
Zbiór $\{r=0\}$ ma miarę $P_Y$ równą zero, więc nie
wpływa na wzory. Funkcja $\beta_i(Y)$ jest
$\sigma(Y)$-mierzalna i ograniczona. Jest zatem
warunkową wartością oczekiwaną funkcji charakterystycznej
$\chi_{\{\Theta=i\}}$ względem $\sigma(Y)$.
Równość sumy wynika z $\sum_i\pi_ip_i=r$ na $\{r>0\}$.
\end{proof}

\begin{przyklad}[Dwie obserwacje Bernoulliego]
Niech $\Theta\in\{0,1\}$ z prawdopodobieństwami
$\mathbb P(\Theta=1)=\pi$ i $\mathbb P(\Theta=0)=1-\pi$.
Warunkowo względem $\Theta=i$ zaobserwujmy
niezależną parę $(Y_1,Y_2)$ zmiennych Bernoulliego
o parametrze $\theta_i\in(0,1)$. Dla obserwacji
$(Y_1,Y_2)=(1,0)$ wzór Bayesa daje
\[
\mathbb P(\Theta=1\mid Y_1=1,Y_2=0)
=\frac{\pi\theta_1(1-\theta_1)}
{\pi\theta_1(1-\theta_1)+(1-\pi)\theta_0(1-\theta_0)}.
\]
Nie potrzeba tu miary Lebesgue'a: wspólną miarą
dominującą jest miara licząca na czteroelementowym zbiorze.
\end{przyklad}

\section{Warunkowa wartość oczekiwana i nierówności}
Warunkową wartość oczekiwaną zdefiniowaliśmy w
rozdziale~\ref{ch:radon-nikodym-18}. Teraz dowiedziemy
jej własności przydatnych w analizie oraz prawdopodobieństwie.

\begin{lemat}[Wyciąganie czynnika mierzalnego]
\label{lem:conditional-pullout-19}
Niech $(\Omega,\mathcal F,\mathbb P)$ będzie przestrzenią
probabilistyczną, $\mathcal G\subseteq\mathcal F$,
$Y\in L^1(\mathbb P)$ oraz $H$ ograniczoną rzeczywistą
funkcją $\mathcal G$-mierzalną. Wówczas
\[
 \mathbb E[HY\mid\mathcal G]
 =H\mathbb E[Y\mid\mathcal G]
 \quad\mathbb P\text{-prawie wszędzie}.
\]
\end{lemat}
\begin{proof}
Dla $H=\chi_B$, $B\in\mathcal G$, tożsamość wynika
z warunku definiującego warunkową wartość oczekiwaną,
ponieważ dla dowolnego $A\in\mathcal G$ całkujemy po
$A\cap B$. Z liniowości wynika przypadek prostych $H$.
Dla ograniczonego $H$ istnieją proste $H_n$, mierzalne
względem $\mathcal G$, takie, że $H_n\to H$ punktowo
i $|H_n|\le\|H\|_\infty+1$. Przez zbieżność zmajoryzowaną
$H_nY\to HY$ i $H_n\mathbb E[Y\mid\mathcal G]
\to H\mathbb E[Y\mid\mathcal G]$ w $L^1$.
Całkując po każdym $A\in\mathcal G$, przechodzimy
do granicy; jednoznaczność daje tezę.
\end{proof}

\begin{twierdzenie}[Warunkowa nierówność Jensena]
\label{thm:conditional-jensen-19}
Niech $(\Omega,\mathcal F,\mathbb P)$ będzie przestrzenią
probabilistyczną, $\mathcal G\subseteq\mathcal F$,
$Y\in L^1(\mathbb P;\mathbb R)$, a
$\Phi:\mathbb R\to\mathbb R$ funkcją wypukłą, wszędzie
skończoną. Załóżmy ponadto, że $\Phi(Y)\in L^1(\mathbb P)$.
Wówczas
\[
 \Phi\bigl(\mathbb E[Y\mid\mathcal G]\bigr)
 \le\mathbb E[\Phi(Y)\mid\mathcal G]
 \quad\mathbb P\text{-prawie wszędzie}.
\]
\end{twierdzenie}
\begin{proof}
Dla każdego $t\in\mathbb Q$ wybieramy liczbę
$a_t\in\partial\Phi(t)$, gdzie
$\partial\Phi(t)$ oznacza zbiór nachyleń prostych
podpierających wykres $\Phi$ w punkcie $t$.
Istnienie takiego nachylenia wynika z monotoniczności
ilorazów różnicowych funkcji wypukłej; możemy np. przyjąć
prawostronną pochodną $\Phi'_+(t)$, która jest skończona.
Wówczas
$\Phi(x)\ge a_tx+\Phi(t)-a_tt$ dla każdego $x\in\mathbb R$.
Różnica
$\Phi(Y)-a_tY-\Phi(t)+a_tt$ jest nieujemna,
więc dodatniość i liniowość warunkowej wartości oczekiwanej
dają
\[
\mathbb E[\Phi(Y)\mid\mathcal G]
\ge a_t\mathbb E[Y\mid\mathcal G]+\Phi(t)-a_tt.
\]
Ponieważ zbiór liczb wymiernych jest przeliczalny,
można wybrać jeden wspólny zbiór pełnej miary, na którym
wszystkie nierówności obowiązują jednocześnie.
Ciągłość $\Phi$ i jej lokalna lipschitzowskość
wynikają z wypukłości i skończoności na całej prostej.
Pozwala to przejść z $t\in\mathbb Q$ do dowolnego $t$:
$\Phi(x)=\sup_{t\in\mathbb Q}
(a_tx+\Phi(t)-a_tt)$.
Podstawiając $x=\mathbb E[Y\mid\mathcal G]$
i biorąc supremum otrzymujemy tezę.
\end{proof}

\begin{wniosek}[Kontrakcja w $L^p$]
\label{cor:conditional-lp-19}
Jeśli $1\le p<\infty$, $Y\in L^p(\mathbb P;\mathbb R)$,
to $\mathbb E[Y\mid\mathcal G]\in L^p(\mathbb P)$ oraz
$\|\mathbb E[Y\mid\mathcal G]\|_p\le\|Y\|_p$.
Dla $p=\infty$ zachodzi analogiczne oszacowanie
normy istotnego supremum.
\end{wniosek}
\begin{proof}
Dla $1\le p<\infty$ stosujemy twierdzenie
\ref{thm:conditional-jensen-19} do $\Phi(t)=|t|^p$.
Całkując otrzymaną nierówność względem $\mathbb P$
i korzystając z własności zachowania całki warunkowej,
dostajemy oszacowanie normy. Jeżeli $|Y|\le M$ prawie
wszędzie, to z monotoniczności
$-M\le\mathbb E[Y\mid\mathcal G]\le M$,
co daje przypadek $p=\infty$.
\end{proof}

\section{Informacja ilościowa: entropia względna}
Dla prawdopodobieństw $Q\ll P$ określamy
\emph{entropię względną} (dywergencję Kullbacka--Leiblera)
\[
 D(Q\Vert P):=\int_X f(x)\log f(x)\,\mathrm dP(x),
 \qquad f=\frac{\mathrm dQ}{\mathrm dP},
\]
przy konwencji $0\log0=0$. Całka ma wartości
w $[0,\infty]$: z $t\log t\ge -e^{-1}$
wynika, że dla prawdopodobieństwa $P$ część ujemna
całkowanej funkcji jest całkowalna. Gdy $Q\not\ll P$,
przyjmujemy $D(Q\Vert P)=+\infty$.

\begin{twierdzenie}[Nieujemność entropii względnej]
\label{thm:gibbs-19}
Dla dowolnych prawdopodobieństw $P,Q$ na $(X,\mathcal F)$
zachodzi $D(Q\Vert P)\ge0$. Jeżeli $Q\ll P$, to
$D(Q\Vert P)=0$ wtedy i tylko wtedy, gdy $Q=P$.
\end{twierdzenie}
\begin{proof}
Dla $t\ge0$ funkcja $\psi(t)=t\log t-t+1$
(z $\psi(0)=1$) jest nieujemna i zeruje się wyłącznie
w $t=1$: na $(0,\infty)$ pochodna jest równa $\log t$,
więc minimum jest w $1$. Jeśli $Q\ll P$ i $f=\mathrm dQ/\mathrm dP$,
to $\int_Xf\,\mathrm dP=1$, a zatem
\[
 D(Q\Vert P)=\int_X\psi(f(x))\,\mathrm dP(x)\ge0.
\]
Jeśli całka jest zerowa, to $\psi(f)=0$ $P$-prawie wszędzie,
a więc $f=1$ $P$-prawie wszędzie, czyli $Q=P$.
W pozostałym przypadku definicja daje nieskończoność.
\end{proof}

\begin{twierdzenie}[Zasada przetwarzania danych dla entropii]
\label{thm:kl-data-processing-19}
Niech $P,Q$ będą prawdopodobieństwami na $(X,\mathcal F)$,
$Q\ll P$, a $T:(X,\mathcal F)\to(Y,\mathcal G)$
odwzorowaniem mierzalnym. Wówczas
\[
 D(T_*Q\Vert T_*P)\le D(Q\Vert P).
\]
\end{twierdzenie}
\begin{proof}
Niech $f=\mathrm dQ/\mathrm dP$, a
$g=\mathrm d(T_*Q)/\mathrm d(T_*P)$.
Z twierdzenia~\ref{thm:image-ac-19} wynika
$g(T)=\mathbb E_P[f\mid\sigma(T)]$ prawie wszędzie.
Dla funkcji wypukłej
$\psi(t)=t\log t-t+1$ na $[0,\infty)$
warunkowa nierówność Jensena pozostaje prawdziwa
również wtedy, gdy $\psi(f)$ nie jest całkowalna:
stosujemy ją do funkcji nieujemnej $\psi$ po odjęciu
wspierającej funkcji afinicznej, lub aproksymujemy
$\psi$ supremum przeliczalnej rodziny prostych
podpierających. Wówczas
\[
 \psi(g(T))\le\mathbb E_P[\psi(f)\mid\sigma(T)]
\]
w rozumieniu rozszerzonej warunkowej wartości oczekiwanej
funkcji nieujemnej. Jeżeli $D(Q\Vert P)=\infty$,
nierówność jest natychmiastowa. W przypadku skończonym
$\psi(f)\in L^1(P)$, a powyższą nierówność można
uzyskać wprost z dowodu twierdzenia
\ref{thm:conditional-jensen-19}, z rozszerzeniem na funkcję
$\psi$ na $[0,\infty)$; nieujemność $f$ i $g$
pozwala używać wyłącznie podpór w punktach dodatnich.
Całkując względem $P$ i stosując definicję miary obrazowej,
otrzymujemy
\[
 D(T_*Q\Vert T_*P)
=\int_Y\psi(g(y))\,\mathrm d(T_*P)(y)
\le\int_X\psi(f(x))\,\mathrm dP(x)=D(Q\Vert P).
\]
\end{proof}

\begin{uwaga}[Interpretacja]
Wzory dla wariacji całkowitej i entropii względnej
pokazują dwa różne sensy zasady, że samo przekształcenie
obserwacji nie może zwiększyć informacji o tym,
z którego z dwóch rozkładów pochodzi wynik. Są to
nierówności matematyczne; nie wymagają aparatu teorii
informacji ani wcześniejszej znajomości statystyki.
\end{uwaga}

\section{Ograniczenia i dalsze kierunki}
Pochodna Radona--Nikodyma jest określona prawie wszędzie,
a nie punktowo kanonicznie. Wartość $f(x)$ w danym punkcie
nie musi być znacząca; w zastosowaniach liczą się całki,
zbiory poziomicowe z dokładnością do zbiorów zerowych
oraz odpowiednie klasy funkcji. Warunkowe rozkłady
jako rodziny miar $P(\,\cdot\mid Y=y)$ wymagają dodatkowych
założeń, np. standardowej borelowskości przestrzeni,
czym zajmiemy się w module o dezintegracji.

\begin{przyklad}[Brak punktowej jednoznaczności rozkładu warunkowego]
Niech $\Omega=(0,1)$, $\mathcal F=\mathcal B((0,1))$,
$\mathbb P=\lambda_1|_{(0,1)}$, a $Y(\omega)=\Theta(\omega)=\omega$.
Dla $y\in(0,1)$ i $A\in\mathcal B((0,1))$ zdefiniujmy
$K(y,A)=\chi_A(y)$. Dla każdego $y$ odwzorowanie
$A\mapsto K(y,A)$ jest miarą Diraca $\delta_y$,
zaś dla każdego $A$ funkcja $y\mapsto K(y,A)$ jest
mierzalna. Ponadto dla dowolnych zbiorów borelowskich
$A,B\subseteq(0,1)$ zachodzi
\[
 \int_B K(y,A)\,\mathrm dP_Y(y)
 =\lambda_1(A\cap B)
 =\mathbb P(\Theta\in A,\,Y\in B).
\]
Zatem $K$ jest rozkładem warunkowym zmiennej
$\Theta$ względem obserwacji $Y$.

Ustalmy teraz $y_0\in(0,1)$ i miarę
probabilistyczną $\rho\ne\delta_{y_0}$ na $(0,1)$. Jeżeli zastąpimy
$K(y_0,\cdot)$ przez $\rho$, pozostawiając
$K(y,\cdot)=\delta_y$ dla $y\ne y_0$, otrzymamy
inne jądro mierzalne spełniające tę samą tożsamość
całkową: zbiór $\{y_0\}$ ma miarę $P_Y$ równą zero.
Warunkowy rozkład jest więc określony jedynie
$P_Y$-prawie wszędzie, nawet w tak prostej przestrzeni
standardowej borelowskiej.
\end{przyklad}
